\documentclass[11pt,a4paper]{article}

% ── Packages ──────────────────────────────────────────────────────────────────
\usepackage[margin=1in]{geometry}
\usepackage{graphicx}
\usepackage{booktabs}
\usepackage{amsmath}
\usepackage{amssymb}
\usepackage{hyperref}
\usepackage{float}
\usepackage{caption}
\usepackage{subcaption}
\usepackage{multirow}
\usepackage{array}
\usepackage{xcolor}
\usepackage{microtype}
\usepackage{parskip}
\usepackage[numbers,sort&compress]{natbib}

\hypersetup{
  colorlinks=true,
  linkcolor=blue!60!black,
  citecolor=blue!60!black,
  urlcolor=blue!60!black
}

% ── Title ─────────────────────────────────────────────────────────────────────
\title{\textbf{Reproducibility Study: Probabilistic Plan Recognition\\
Using Off-the-Shelf Classical Planners}\\[6pt]
\large COMS7044A --- Reproducibility Assignment}
\author{Sibonelo Mlambo 536187}
\date{20 May 2026}

% ── Document ──────────────────────────────────────────────────────────────────
\begin{document}

\maketitle
\tableofcontents
\clearpage

% ══════════════════════════════════════════════════════════════════════════════
\section*{Reproducibility Summary}
% ══════════════════════════════════════════════════════════════════════════════

This report presents a reproducibility study of the probabilistic plan-recognition
method proposed by Ram\'{i}rez and Geffner~\cite{ramirez2010} (AAAI 2010).
The method infers an agent's most-likely goal from a sequence of partial
observations by comparing the cost of observation-compliant versus
non-compliant plans and combining the resulting likelihood with a prior through
Bayes' rule.

The implementation covers three planning domains---\textbf{Blocksworld},
\textbf{Logistics}, and \textbf{Campus}.  All experiments were repeated with
50~independent trials per condition and fixed random seeds.  Results are
reported as means with standard deviations and 95\,\% confidence intervals.
The main underspecified quantity in the original paper---the inverse-temperature
parameter $\beta$---is treated as a sensitivity variable and swept across
$\beta \in \{0.1, 0.25, 0.5, 1, 2, 5, 10\}$.

% ══════════════════════════════════════════════════════════════════════════════
\section{Introduction}
% ══════════════════════════════════════════════════════════════════════════════

Goal recognition is the problem of inferring which goal an observed agent is
most likely pursuing from a (possibly incomplete) sequence of its actions.
Ram\'{i}rez and Geffner reduce this to repeated classical planning: for each
candidate goal $G$ they compute the cheapest plan that \emph{complies} with
the observations and the cheapest plan that \emph{does not comply}, take the
difference in costs $\Delta(G,O)$, and exponentiate it through a Boltzmann
function to obtain a likelihood.  A Bayesian posterior over goals is then used
to rank candidates.

The method is elegant and computationally practical because it treats the
planner as a black box.  This study assesses whether the core claims of the
paper are reproducible in a clean Python re-implementation that uses a
built-in A$^*$ planner (with an optional Fast~Downward interface) and evaluates
the robustness of the approach under four observation-noise models.

% ══════════════════════════════════════════════════════════════════════════════
\section{Scope of Reproducibility}
% ══════════════════════════════════════════════════════════════════════════════

The following testable claims are evaluated:

\begin{enumerate}
  \item \textbf{Observation completeness.} Increasing the percentage of observed
        actions improves top-1 goal-recognition accuracy.
  \item \textbf{Posterior concentration.} The true goal receives a higher
        posterior probability when more observations are available.
  \item \textbf{Beta sensitivity.} The parameter $\beta$ controls the sharpness
        of the posterior; larger $\beta$ concentrates mass on the most
        likely goal.
  \item \textbf{Noise robustness.} Observation noise (missing, corrupt, spurious,
        or mixed actions) reduces recognition accuracy, with the
        degree of degradation depending on the noise type and rate.
  \item \textbf{Generality.} The planner-cost approach is reproducible across
        multiple PDDL-style planning domains.
\end{enumerate}

% ══════════════════════════════════════════════════════════════════════════════
\section{Background and Method}
% ══════════════════════════════════════════════════════════════════════════════

\subsection{Probabilistic Plan Recognition}

Let $\mathcal{G}$ be the set of candidate goals and $O = \langle o_1,\ldots,o_k\rangle$
a sequence of observed actions.  For each goal $G \in \mathcal{G}$, define

\begin{align}
  \Delta(G, O) &= \operatorname{cost}(\Pi^-(G,O)) - \operatorname{cost}(\Pi^+(G,O)),
\end{align}

\noindent where $\Pi^+(G,O)$ is an optimal plan for $G$ that \emph{includes} all
observed actions and $\Pi^-(G,O)$ is an optimal plan that is \emph{not} required
to do so.  The Boltzmann likelihood is

\begin{equation}
  P(O \mid G) \;\propto\; e^{\,\beta\,\Delta(G,O)},
\end{equation}

\noindent and the posterior (assuming a uniform prior) is

\begin{equation}
  P(G \mid O) = \frac{e^{\,\beta\,\Delta(G,O)}}{\displaystyle\sum_{G' \in \mathcal{G}} e^{\,\beta\,\Delta(G',O)}}.
\end{equation}

The top-1 predicted goal is $\hat{G} = \arg\max_{G} P(G \mid O)$.

\subsection{Ambiguities}

The original paper does not specify the value of $\beta$ used in its
experiments, which is the primary source of irreproducibility.  This study
therefore treats $\beta$ as a free parameter and reports results for
$\beta \in \{0.1, 0.25, 0.5, 1, 2, 5, 10\}$.

% ══════════════════════════════════════════════════════════════════════════════
\section{Methodology}
% ══════════════════════════════════════════════════════════════════════════════

\subsection{Implementation}

The pipeline is written in Python.  Each domain is self-contained and exposes:

\begin{itemize}
  \item PDDL-style domain and problem files.
  \item An observation loader that samples action-prefix traces at configurable
        observation percentages (25\,\%, 50\,\%, 75\,\%, 100\,\%).
  \item A planner interface backed by a domain-specific A$^*$ implementation;
        Fast~Downward can be substituted via a command-line flag.
  \item A Bayesian scoring layer (\texttt{BayesianGoalScorer}) implementing
        Equations~(1)--(3).
  \item An experiment runner that sweeps $\beta$, observation percentage, noise
        mode, and noise rate across 50 independent trials per condition.
  \item A plotting module that writes all figures used in this report.
\end{itemize}

\subsection{Domains}

\paragraph{Blocksworld.} 
\begin{itemize}
    \item A 3-block stacking domain.  \item Candidate goals are all
permutations of valid stack configurations.  
\item Plans are computed with a
domain-specific A$^*$ planner. 
\end{itemize} 

\paragraph{Logistics.}  
\begin{itemize}
\item A package-delivery domain with trucks and locations.
\item The increased branching factor makes this domain harder; fewer observations
are sufficient to distinguish goals.
\end{itemize}


\paragraph{Campus.} 
\begin{itemize}
\item A navigation domain over a campus graph.  
\item Goals correspond
to destination rooms.
\item Paths are found with A$^*$ over the campus graph.
\end{itemize}


\subsection{Experimental Design}

\begin{itemize}
  \item \textbf{$\beta$ values:} $\{0.1, 0.25, 0.5, 1, 2, 5, 10\}$.
  \item \textbf{Observation percentages:} $\{25\,\%, 50\,\%, 75\,\%, 100\,\%\}$.
  \item \textbf{Noise modes:} none, missing, corrupt, spurious, mixed.
  \item \textbf{Noise rates:} $\{0, 0.1, 0.2, 0.3\}$.
  \item \textbf{Trials per condition:} 50, with a fixed random seed.
\end{itemize}

\paragraph{Metrics.}  Top-1 accuracy, mean true-goal posterior probability,
and mean true-goal rank, each reported with standard deviation and 95\,\%
confidence interval.

% ══════════════════════════════════════════════════════════════════════════════
\section{Results}
% ══════════════════════════════════════════════════════════════════════════════

\subsection{Clean Observations (No Noise)}

\subsubsection{Top-1 Accuracy}

Table~\ref{tab:clean_accuracy} reports top-1 accuracy at $\beta = 1$ for all
three domains across observation percentages.

\begin{table}[H]
\centering
\caption{Top-1 accuracy (mean $\pm$ 95\,\% CI half-width) at $\beta = 1.0$,
         clean observations.}
\label{tab:clean_accuracy}
\begin{tabular}{lccc}
\toprule
\textbf{Obs.\ \%} & \textbf{Blocksworld} & \textbf{Logistics} & \textbf{Campus} \\
\midrule
25\,\%  & $0.64 \pm 0.133$ & $0.26 \pm 0.122$ & $0.20 \pm 0.111$ \\
50\,\%  & $1.00 \pm 0.000$ & $0.26 \pm 0.122$ & $0.20 \pm 0.111$ \\
75\,\%  & $1.00 \pm 0.000$ & $0.26 \pm 0.122$ & $0.78 \pm 0.115$ \\
100\,\% & $1.00 \pm 0.000$ & $0.46 \pm 0.138$ & $1.00 \pm 0.000$ \\
\bottomrule
\end{tabular}
\end{table}

Blocksworld achieves perfect accuracy from 50\,\% observations onward.
Campus requires 100\,\% for perfect accuracy.  Logistics plateaus at
46\,\% even with complete observations, suggesting ambiguous cost
structure in the Logistics domain.

\subsubsection{Effect of $\beta$ on Accuracy and Posterior Probability}

Figures~\ref{fig:bw_clean}--\ref{fig:campus_clean} show accuracy and
true-goal posterior probability as $\beta$ varies for each domain.

\begin{figure}[H]
  \centering
  \begin{subfigure}[b]{0.48\textwidth}
    \includegraphics[width=\textwidth]{figures/blocksworld/clean_accuracy_vs_beta.png}
    \caption{Accuracy vs.\ $\beta$}
  \end{subfigure}\hfill
  \begin{subfigure}[b]{0.48\textwidth}
    \includegraphics[width=\textwidth]{figures/blocksworld/clean_true_goal_probability_vs_beta.png}
    \caption{True-goal posterior vs.\ $\beta$}
  \end{subfigure}
  \caption{Blocksworld --- clean observations.}
  \label{fig:bw_clean}
\end{figure}

\begin{figure}[H]
  \centering
  \begin{subfigure}[b]{0.48\textwidth}
    \includegraphics[width=\textwidth]{figures/blocksworld/clean_true_goal_rank_vs_beta.png}
    \caption{True-goal rank vs.\ $\beta$}
  \end{subfigure}
  \caption{Blocksworld --- true-goal rank under clean observations.}
  \label{fig:bw_rank}
\end{figure}

\begin{figure}[H]
  \centering
  \begin{subfigure}[b]{0.48\textwidth}
    \includegraphics[width=\textwidth]{figures/logistics/clean_accuracy_vs_beta.png}
    \caption{Accuracy vs.\ $\beta$}
  \end{subfigure}\hfill
  \begin{subfigure}[b]{0.48\textwidth}
    \includegraphics[width=\textwidth]{figures/logistics/clean_true_goal_probability_vs_beta.png}
    \caption{True-goal posterior vs.\ $\beta$}
  \end{subfigure}
  \caption{Logistics --- clean observations.}
  \label{fig:log_clean}
\end{figure}

\begin{figure}[H]
  \centering
  \begin{subfigure}[b]{0.48\textwidth}
    \includegraphics[width=\textwidth]{figures/campus/clean_accuracy_vs_beta.png}
    \caption{Accuracy vs.\ $\beta$}
  \end{subfigure}\hfill
  \begin{subfigure}[b]{0.48\textwidth}
    \includegraphics[width=\textwidth]{figures/campus/clean_true_goal_probability_vs_beta.png}
    \caption{True-goal posterior vs.\ $\beta$}
  \end{subfigure}
  \caption{Campus --- clean observations.}
  \label{fig:campus_clean}
\end{figure}

In all three domains, larger $\beta$ monotonically increases the
true-goal posterior probability, consistent with the theoretical role
of $\beta$ as an inverse temperature.  For Blocksworld, accuracy is
insensitive to $\beta$ once $\geq 50\,\%$ of observations are
available, because the cost difference $\Delta(G,O)$ unambiguously
ranks the true goal first for any positive $\beta$.

\subsubsection{Posterior Probability at $\beta = 1.0$}

Table~\ref{tab:clean_prob} shows the mean true-goal posterior probability
at $\beta = 1.0$.

\begin{table}[H]
\centering
\caption{Mean true-goal posterior probability at $\beta = 1.0$, clean observations.}
\label{tab:clean_prob}
\begin{tabular}{lccc}
\toprule
\textbf{Obs.\ \%} & \textbf{Blocksworld} & \textbf{Logistics} & \textbf{Campus} \\
\midrule
25\,\%  & 0.393 & 0.135 & 0.200 \\
50\,\%  & 0.827 & 0.075 & 0.289 \\
75\,\%  & 0.973 & 0.000 & 0.429 \\
100\,\% & 0.996 & 0.619 & 0.592 \\
\bottomrule
\end{tabular}
\end{table}

Blocksworld shows a clear monotone increase in posterior probability
with observation percentage.  Campus shows the same trend.  Logistics
behaves anomalously: the posterior falls to zero at 75\,\% before
recovering at 100\,\%, indicating that intermediate observations can be
misleading in domains with symmetrical cost structure.

% ─────────────────────────────────────────────────────────────────────────────
\subsection{Noisy Observations}
% ─────────────────────────────────────────────────────────────────────────────

Four noise models were evaluated at rates $r \in \{0.1, 0.2, 0.3\}$:

\begin{description}
  \item[Missing.] Actions are randomly dropped from the observation
    sequence with probability $r$.
  \item[Corrupt.] Actions are replaced with random valid actions with
    probability $r$.
  \item[Spurious.] Extra random actions are inserted with probability $r$.
  \item[Mixed.] A combination of missing, corrupt, and spurious noise.
\end{description}

\subsubsection{Accuracy at $\beta = 1.0$, $100\,\%$ Observations}

Table~\ref{tab:noise_accuracy} reports accuracy under each noise mode
and rate at $\beta = 1.0$ with full observations.

\begin{table}[H]
\centering
\caption{Top-1 accuracy under observation noise ($\beta = 1.0$, $100\,\%$ observations).}
\label{tab:noise_accuracy}
\setlength{\tabcolsep}{5pt}
\begin{tabular}{llccc}
\toprule
\textbf{Noise mode} & \textbf{Domain} & \textbf{rate = 0.1} & \textbf{rate = 0.2} & \textbf{rate = 0.3} \\
\midrule
\multirow{3}{*}{Missing}
  & Blocksworld & 0.88 & 0.76 & 0.68 \\
  & Logistics   & 0.32 & 0.26 & 0.28 \\
  & Campus      & 0.82 & 0.60 & 0.52 \\
\midrule
\multirow{3}{*}{Corrupt}
  & Blocksworld & 0.74 & 0.60 & 0.46 \\
  & Logistics   & 0.26 & 0.22 & 0.24 \\
  & Campus      & 0.84 & 0.76 & 0.64 \\
\midrule
\multirow{3}{*}{Spurious}
  & Blocksworld & 0.88 & 0.78 & 0.66 \\
  & Logistics   & 0.28 & 0.36 & 0.40 \\
  & Campus      & 0.88 & 0.82 & 0.72 \\
\midrule
\multirow{3}{*}{Mixed}
  & Blocksworld & 0.62 & 0.48 & 0.32 \\
  & Logistics   & 0.26 & 0.28 & 0.28 \\
  & Campus      & 0.66 & 0.56 & 0.50 \\
\bottomrule
\end{tabular}
\end{table}

Several patterns are evident.  Blocksworld degrades gracefully under
missing and spurious noise but drops sharply under mixed noise.
Logistics accuracy is near-constant across all conditions, reflecting
its baseline difficulty even without noise.  Campus is the most
robust domain, retaining $\geq 0.50$ accuracy even at 30\,\% mixed noise.

\subsubsection{Noise Heatmaps}

Figures~\ref{fig:bw_noise}--\ref{fig:campus_noise} display
accuracy heatmaps over noise rate and $\beta$ for each noise mode
and domain.

\begin{figure}[H]
  \centering
  \begin{subfigure}[b]{0.48\textwidth}
    \includegraphics[width=\textwidth]{figures/blocksworld/heatmap_missing_noise_accuracy.png}
    \caption{Missing noise}
  \end{subfigure}\hfill
  \begin{subfigure}[b]{0.48\textwidth}
    \includegraphics[width=\textwidth]{figures/blocksworld/heatmap_corrupt_noise_accuracy.png}
    \caption{Corrupt noise}
  \end{subfigure}
  \\[6pt]
  \begin{subfigure}[b]{0.48\textwidth}
    \includegraphics[width=\textwidth]{figures/blocksworld/heatmap_spurious_noise_accuracy.png}
    \caption{Spurious noise}
  \end{subfigure}\hfill
  \begin{subfigure}[b]{0.48\textwidth}
    \includegraphics[width=\textwidth]{figures/blocksworld/heatmap_mixed_noise_accuracy.png}
    \caption{Mixed noise}
  \end{subfigure}
  \caption{Blocksworld --- accuracy heatmaps (noise rate vs.\ $\beta$).}
  \label{fig:bw_noise}
\end{figure}

\begin{figure}[H]
  \centering
  \begin{subfigure}[b]{0.48\textwidth}
    \includegraphics[width=\textwidth]{figures/logistics/heatmap_missing_noise_accuracy.png}
    \caption{Missing noise}
  \end{subfigure}\hfill
  \begin{subfigure}[b]{0.48\textwidth}
    \includegraphics[width=\textwidth]{figures/logistics/heatmap_corrupt_noise_accuracy.png}
    \caption{Corrupt noise}
  \end{subfigure}
  \\[6pt]
  \begin{subfigure}[b]{0.48\textwidth}
    \includegraphics[width=\textwidth]{figures/logistics/heatmap_spurious_noise_accuracy.png}
    \caption{Spurious noise}
  \end{subfigure}\hfill
  \begin{subfigure}[b]{0.48\textwidth}
    \includegraphics[width=\textwidth]{figures/logistics/heatmap_mixed_noise_accuracy.png}
    \caption{Mixed noise}
  \end{subfigure}
  \caption{Logistics --- accuracy heatmaps (noise rate vs.\ $\beta$).}
  \label{fig:log_noise}
\end{figure}

\begin{figure}[H]
  \centering
  \begin{subfigure}[b]{0.48\textwidth}
    \includegraphics[width=\textwidth]{figures/campus/heatmap_missing_noise_accuracy.png}
    \caption{Missing noise}
  \end{subfigure}\hfill
  \begin{subfigure}[b]{0.48\textwidth}
    \includegraphics[width=\textwidth]{figures/campus/heatmap_corrupt_noise_accuracy.png}
    \caption{Corrupt noise}
  \end{subfigure}
  \\[6pt]
  \begin{subfigure}[b]{0.48\textwidth}
    \includegraphics[width=\textwidth]{figures/campus/heatmap_spurious_noise_accuracy.png}
    \caption{Spurious noise}
  \end{subfigure}\hfill
  \begin{subfigure}[b]{0.48\textwidth}
    \includegraphics[width=\textwidth]{figures/campus/heatmap_mixed_noise_accuracy.png}
    \caption{Mixed noise}
  \end{subfigure}
  \caption{Campus --- accuracy heatmaps (noise rate vs.\ $\beta$).}
  \label{fig:campus_noise}
\end{figure}

% ══════════════════════════════════════════════════════════════════════════════
\section{Discussion}
% ══════════════════════════════════════════════════════════════════════════════

\subsection{Claim-by-Claim Assessment}

\paragraph{Claim 1: More observations improve accuracy.}
\textbf{Supported} for Blocksworld (64\,\% $\to$ 100\,\%) and Campus
(20\,\% $\to$ 100\,\%).  \textbf{Partially supported} for Logistics, which
shows a non-monotone profile (accuracy dips at 75\,\% before recovering
at 100\,\%).  This is attributable to cost symmetry in intermediate
plan prefixes.

\paragraph{Claim 2: Posterior probability concentrates on the true goal.}
\textbf{Supported} in Blocksworld and Campus.  In Logistics the posterior
even reaches zero at 75\,\% observations, indicating the planner assigns
identical compliant and non-compliant costs at that point---a limitation
noted but not addressed in the original paper.

\paragraph{Claim 3: $\beta$ controls posterior sharpness.}
\textbf{Strongly supported.}  In all three domains, increasing $\beta$
monotonically raises the true-goal posterior probability.  Accuracy is
less sensitive because rank-1 is maintained once the cost gap is positive,
regardless of magnitude.

\paragraph{Claim 4: Noise reduces accuracy.}
\textbf{Supported.}  All noise types reduce accuracy relative to the
clean baseline, with mixed noise causing the steepest degradation.
Spurious noise is least damaging in Blocksworld and Campus; corrupt
noise is most damaging.

\paragraph{Claim 5: The method generalises across domains.}
\textbf{Partially supported.}  The approach works well on Blocksworld and
Campus.  Logistics presents persistent challenges likely due to cost
ties in the domain's action structure.

\subsection{Reproducibility Limitations}

\begin{enumerate}
  \item \textbf{Unspecified $\beta$.}  The original paper gives no value for
        $\beta$, making exact numerical comparison impossible.  This study
        uses $\beta = 1$ as the reference and sweeps a wide range.
  \item \textbf{Planner differences.}  The built-in A$^*$ planner may return
        different optimal plans from Fast~Downward when there are ties,
        affecting $\Delta(G,O)$.
  \item \textbf{Benchmark generation.}  The original benchmarks are not
        publicly available; this study uses manually constructed problem
        instances of comparable size.
  \item \textbf{Domain scale.}  All domains are small (3 blocks, 5 goals per
        domain) to permit exhaustive repeated trials; the original paper
        uses larger instances.
\end{enumerate}

% ══════════════════════════════════════════════════════════════════════════════
\section{Conclusion}
% ══════════════════════════════════════════════════════════════════════════════

This study provides a clean, self-contained Python reproduction of the
probabilistic plan-recognition framework of Ram\'{i}rez and Geffner.
The central idea---using planner-derived cost differences as a proxy for
observational likelihood---is confirmed to work robustly on two out of
three domains under a range of observation percentages, $\beta$ values,
and noise conditions.  The Logistics domain reveals a domain-specific
failure mode related to cost symmetry in intermediate plan prefixes.

The main obstacle to exact numerical reproduction is the unspecified
$\beta$ parameter; this study addresses this by reporting a full
sensitivity sweep.  The implementation, data, and plots are packaged in
the accompanying repository for independent verification.

% ── References ────────────────────────────────────────────────────────────────
\bibliographystyle{plainnat}
\begin{thebibliography}{1}

\bibitem{ramirez2010}
Miquel Ram\'{i}rez and Hector Geffner.
\newblock Probabilistic plan recognition using off-the-shelf classical planners.
\newblock In \emph{Proceedings of the 24th AAAI Conference on Artificial
  Intelligence (AAAI-10)}, pages 1121--1126, 2010.

\end{thebibliography}

\end{document}
